\section{Retain state across helper calls}
The whole-buffer decoder calls validation, which itself calls the classifier. Each
\texttt{jal} replaces \texttt{ra}, so a nonleaf function must save its caller's return
address before making the next call. The original source, destination, count and capacity
also need stable homes. Argument and temporary registers may be overwritten by a helper.

After validation returns, check status before interpreting its output count. Then
compare the original unsigned capacity with that count. In the illustrated frame,
\texttt{s3} retains capacity until this comparison; it becomes the fixed total output
count only after the comparison succeeds.

\Needspace{13\baselineskip}
\MARSCodeLabel{Use a helper result only after its status and capacity checks.}
\begin{minted}[samepage=true]{mars-idea}
    jal b64_validate
    nop
    bnez $v0, decode_failed
    nop
    sltu $t0, $s3, $v1
    bnez $t0, decode_capacity
    nop
    move $s3, $v1
\end{minted}

The error return sets the produced count to zero and restores the same full-width
saved state as the successful return. An early count error or empty request has no
allocated frame and returns directly. Keep those paths distinct; restoring an
unallocated frame would change the caller's stack pointer.

Set nonzero high bits in preserved registers while testing, and place code and data
above 4 GiB. A word restore can leave plausible low bits while losing a saved value
or address. The width of the encoded characters does not determine the width of the
return address, stack pointer or buffer pointers.
